%----------------------------------------------------------------------------
\chapter{Implementation}
\label{sec:Implementation}
%----------------------------------------------------------------------------
\providecommand{\reasontodo}[1]{\textcolor{red}{\textbf{[TODO: #1]}}}

This chapter describes how the design of \autoref{sec:Design} became a working system, from the point of view of the developer. \autoref{sec:ImplEnvironment} presents the environment and the tools, and \autoref{sec:ImplRepos} the repositories. \autoref{sec:ImplBootstrap} and \autoref{sec:ImplGitOps} follow the two layers of the platform, \autoref{sec:ImplNetworking} explains how the platform is reached, and \autoref{sec:ImplObservability} how its metrics are collected. \autoref{sec:ImplDemo} covers the demo application and its release pipeline, and \autoref{sec:ImplMechanisms} the detection mechanisms, with the points where the design could not be implemented directly and the workarounds that were needed. \autoref{sec:ImplProd} closes with the state of the production environment.

%----------------------------------------------------------------------------
\section{Environment and tools}
\label{sec:ImplEnvironment}
%----------------------------------------------------------------------------
The platform runs on two virtual machines with 8\,GB of memory and 4 CPU cores each, both with Ubuntu 24.04, reachable over SSH and connected to the control machine by the NetBird VPN. \texttt{prgdlv-test} runs the development cluster. \texttt{prgdlv-auto} runs Gitea with its container registry and an \texttt{act\_runner} CI runner as host services, and is reserved for the production cluster. The control machine, a laptop on the same VPN, runs ansible-core, \texttt{kubectl}, Helm, SOPS and age, and the Python standard library is the only dependency of the trial runner. Only SSH access with \texttt{sudo} and a running NetBird client are prepared on the VMs by hand; everything else is installed by the bootstrap.

Every component is pinned to an exact version (\autoref{tab:ImplVersions}), in the Ansible requirements and in the Argo CD Applications. An environment rebuilt months later therefore gets the same software as the one it replaces (R6), and the results of the experiments refer to known versions, which turned out to matter: the anomaly analysed in \autoref{sec:EvalDeadline} is a defect of one specific version of Argo Rollouts.

\begin{table}[ht]
	\footnotesize
	\centering
	\begin{tabular}{ l l }
		\toprule
		Component & Version \\
		\midrule
		k3s & v1.36.5+k3s1 \\
		k3s Ansible collection & 1.2.2 \\
		Argo CD (Helm chart) & v3.5.5 (10.10.3) \\
		Argo Rollouts (Helm chart) & v1.10.1 (2.43.7) \\
		Istio (\texttt{base}, \texttt{istiod}, \texttt{gateway} charts) & 1.30.6 \\
		cert-manager & v1.21.2 \\
		Prometheus (Helm chart) & 29.37.0 \\
		Grafana (Helm chart) & 13.2.4 (13.7.0) \\
		Kiali server (Helm chart) & 2.33.0 \\
		\bottomrule
	\end{tabular}
	\caption{Pinned versions of the platform components.}
	\label{tab:ImplVersions}
\end{table}

The memory budget of R5 held. Before the dashboards were added, the development node used 3.7 of its 8\,GiB, with Prometheus as the largest consumer at about 490\,MiB. Grafana settled at about 385\,MiB when idle. All platform components have explicit resource requests and limits and run with a single replica.

%----------------------------------------------------------------------------
\section{Repositories}
\label{sec:ImplRepos}
%----------------------------------------------------------------------------
The work lives in two Git repositories on Gitea. The application repository, \texttt{demo}, contains the demo application, its Dockerfile and its CI workflow. The platform repository, \texttt{rollback-platform}, contains everything else; its layout follows the two layers of the design (\autoref{lst:ImplLayout}). The bootstrap and the GitOps manifests share one repository; I considered separating them but judged it not worth the effort.

\begin{lstlisting}[caption={Layout of the platform repository (abridged).},label=lst:ImplLayout]
ansible/
  site.yml                      # runs the two playbooks in order
  playbooks/01-k3s.yml          # k3s on the VM
  playbooks/02-bootstrap.yml    # secrets, Argo CD, root Application
  inventories/{dev,prod}/       # hosts, variables, SOPS-encrypted secrets
  roles/argocd_bootstrap/
gitops/
  clusters/{dev,prod}/          # one Argo CD Application per component
  platform/                     # Helm values and manifests of the platform
  apps/demo/                    # canary Rollout, analyses, components s0-s4
  apps/demo-bg/                 # blue-green Rollout, load Job, analyses, b0-b3
tools/trial-runner/             # the experiment runner (Evaluation)
docs/                           # DNS, CA certificates, test plans
\end{lstlisting}

%----------------------------------------------------------------------------
\section{Bootstrap with Ansible}
\label{sec:ImplBootstrap}
%----------------------------------------------------------------------------
\texttt{site.yml} runs two playbooks, and one command with the inventory of an environment builds it from a fresh VM: \texttt{ansible-playbook -i inventories/dev/hosts.yml site.yml}.

\paragraph{k3s.} The first playbook runs the official k3s Ansible collection with a single host in the server group and no agents, which creates a single-node cluster. Traefik is disabled with \texttt{--disable traefik}, ServiceLB stays enabled. The cluster token is read from the encrypted secrets, and the cluster's kubeconfig is merged into the control machine's under the context \texttt{k3s-dev} or \texttt{k3s-prod}. Gitea's registry serves only plain HTTP, which the container runtime of k3s refuses by default; a \texttt{registries.yaml} in the inventory declares the registry as an HTTP mirror. k3s reads this file only at start-up, so a small play after the collection restarts k3s when the file is newer than the running process. The images are public in the registry, so no pull secret is needed.

\paragraph{Secrets and Argo CD.} The second playbook runs the role \texttt{argocd\_bootstrap} on the control machine. The encrypted inventory files are decrypted in memory by the SOPS variables plugin of Ansible, so no plaintext secret is ever written to disk. The role creates the namespaces and the three Secrets of \autoref{sec:DesignLayers}, with \texttt{no\_log} so that their values do not appear in Ansible's output, and installs Argo CD with Helm. Argo CD runs with single replicas, without Dex and the notifications controller, which are not used, and in insecure mode, because TLS is terminated at the Istio gateway in front of it. Finally, the role applies the root Application, which points at \texttt{gitops/clusters/<env>/} on the environment's branch.

%----------------------------------------------------------------------------
\section{GitOps layer}
\label{sec:ImplGitOps}
%----------------------------------------------------------------------------
The root Application uses the app-of-apps pattern: the manifests it applies are themselves Applications, one per component~\cite{ArgoCDBootstrap}. Adding a component means adding one file to \texttt{gitops/clusters/<env>/}, and deleting the file removes it, because pruning is enabled. All Applications sync automatically and heal themselves.

\paragraph{Order.} The components depend on each other through their custom resource definitions: a \texttt{ClusterIssuer} can only be created after cert-manager has installed its CRDs, a VirtualService after Istio, and a Rollout after Argo Rollouts. Sync waves order the Applications (\autoref{tab:ImplWaves})~\cite{ArgoCDSyncWaves}. By default, however, Argo CD considers a child Application healthy as soon as it exists, so the waves would not wait for anything. A custom health check, a short Lua script in the Argo CD configuration, makes the health of an Application equal to the health of its resources, and only then does each wave wait for the previous one to become healthy.

\begin{table}[ht]
	\footnotesize
	\centering
	\begin{tabular}{ c p{4.6cm} p{6.8cm} }
		\toprule
		Wave & Applications & Provides \\
		\midrule
		0 & cert-manager & certificate CRDs and controller \\
		1 & istio-base & Istio CRDs \\
		2 & istiod & mesh control plane \\
		3 & istio-ingressgateway, networking & gateway; CA issuer, wildcard certificate, Gateway, UI routes \\
		4 & prometheus, grafana, kiali & metrics and dashboards \\
		5 & argo-rollouts & Rollout CRDs and controller \\
		6 & demo, demo-bg & the applications under test \\
		\bottomrule
	\end{tabular}
	\caption{Sync waves of the platform.}
	\label{tab:ImplWaves}
\end{table}

\paragraph{Charts and values.} Third-party components are installed from their Helm charts. Each such Application has two sources: the chart from its public repository, at a pinned version, and the platform repository, referenced as \texttt{\$values}, which supplies a common values file and one per environment~\cite{ArgoCDMultiSource}. The difference between the environments is thereby reduced to small files such as \texttt{values-dev.yaml}, which sets, for example, a Prometheus retention of 24 hours instead of seven days. The applications of this work are plain manifests, organised as a Kustomize base with an overlay per environment.

\paragraph{Fields that others own.} Two kinds of objects are changed in the cluster by someone other than Argo CD, and Argo CD must not reset them. istiod rewrites its own validating webhook configuration at run time, and Argo Rollouts rewrites the routes of the demo application's VirtualService during every canary. Both are excluded from comparison with \texttt{ignoreDifferences}~\cite{ArgoCDDiffing}, and the sync option \texttt{RespectIgnoreDifferences} excludes them from synchronisation as well. For the VirtualService, the whole route list \texttt{/spec/http} had to be excluded, not only the first route, because mirroring inserts a managed route in front of the weighted one and shifts its index (\autoref{sec:ImplMirror}). The trade-off is that Argo CD no longer restores the routes if someone edits them by hand; it still manages the host names and gateways of the VirtualService. Two further options were needed: server-side apply, because the CRDs of cert-manager and Istio are too large for the annotation that a client-side apply stores, and skipping the dry run for resources whose CRD is still being installed.

%----------------------------------------------------------------------------
\section{Networking and access}
\label{sec:ImplNetworking}
%----------------------------------------------------------------------------
\paragraph{Path of a request.} The Istio ingress gateway is a Service of type \texttt{LoadBalancer}, so ServiceLB exposes it on ports 80 and 443 of the node~\cite{K3sServiceLB}. A request from the VPN to port 443 of the node passes three address translations: the host-port rule of the node forwards it to the ServiceLB pod, that pod forwards it to the cluster IP of the gateway Service, and the Service forwards it to the gateway pod, where Envoy terminates TLS. A host port is a packet-filter rule, not a listening socket, so the ports look free to tools such as \texttt{ss}, and a program on the host listening on the same port would be bypassed silently. Gitea on prgdlv-auto listens on port 3000 and is therefore not affected.

\paragraph{TLS and names.} In \texttt{istio-system}, the networking Application creates the CA issuer from the root CA Secret, a wildcard certificate for \texttt{<env>.internal} and \texttt{*.<env>.internal} (valid for 90 days and renewed 30 days before it expires) and a \texttt{Gateway} that redirects HTTP to HTTPS and serves the certificate. The names are resolved by NetBird's DNS: \texttt{*.dev.internal} points to the development VM and \texttt{*.prod.internal} to the production VM. The public certificates of the two root CAs are kept in the repository and are installed by hand into the trust store of each device that uses the platform. \autoref{tab:ImplUIs} lists the user interfaces, all routed through the same gateway.

\begin{table}[ht]
	\footnotesize
	\centering
	\begin{tabular}{ l l l }
		\toprule
		Host & Service & Access \\
		\midrule
		\texttt{app.dev.internal} & canary demo application & open within the VPN \\
		\texttt{bg.dev.internal} & blue-green demo application & open within the VPN \\
		\texttt{argocd.dev.internal} & Argo CD & administrator login \\
		\texttt{rollouts.dev.internal} & Argo Rollouts dashboard & none (can promote and abort) \\
		\texttt{prometheus.dev.internal} & Prometheus & no login \\
		\texttt{grafana.dev.internal} & Grafana & anonymous viewers \\
		\texttt{kiali.dev.internal} & Kiali & anonymous, read-only \\
		\bottomrule
	\end{tabular}
	\caption{Services exposed by the development environment.}
	\label{tab:ImplUIs}
\end{table}

\paragraph{The VPN drops forwarded traffic.} After the first deployment, the certificate chain was valid and \texttt{https://app.dev.internal} answered on the VM itself, but from the control machine ports 80 and 443 timed out, while the Kubernetes API on port 6443 of the same address worked. The cause lies in where the packets are rewritten. The host-port rule translates the destination to a pod address \emph{before} the kernel's routing decision, so a request to the gateway counts as forwarded traffic, whereas the API server listens on the host and its packets stay local. NetBird installs its own packet filter, which accepts local traffic from permitted peers but drops forwarded traffic arriving from the VPN unless the peer has been declared a routing peer for the destination network. An accept rule in another table, such as the one k3s adds, cannot override a drop in NetBird's table. The default is reasonable: without it, any host that runs Kubernetes or Docker, which enable IP forwarding, would expose its internal networks to every peer.

Three fixes were possible. Disabling NetBird's firewall on the VM would have opened every port of the VM to every peer. Running the gateway in the host network would avoid the translation altogether, but would bind Envoy directly to the host's ports. The chosen fix declares the pod network of the cluster, \texttt{10.42.0.0/16}, as a NetBird network with the VM as its routing peer, and allows the control machine to reach it~\cite{NetBirdNetworks}. NetBird then marks these packets as permitted, and all other forwarded traffic stays blocked. The production VM will need the same route.

%----------------------------------------------------------------------------
\section{Observability}
\label{sec:ImplObservability}
%----------------------------------------------------------------------------
\paragraph{How Prometheus finds the metrics.} Prometheus discovers its targets through the Kubernetes API and selects them with relabelling rules~\cite{PromConfig}. Most scrape jobs are the defaults of the Prometheus chart. I added one job, \texttt{envoy-stats}: it scrapes every Istio proxy on its metrics port 15090, recognised by the name of the container port, and copies all labels of the pod onto the scraped series. This copy is what lets the checks select the new version by the label \texttt{rollouts\_pod\_template\_hash} (\autoref{sec:DesignChecks}). Memory is read from the kubelet's built-in cAdvisor, which measures every container from its cgroup, so the application exposes no memory metric of its own. The metrics of the Argo Rollouts controller are scraped through an annotation on its metrics Service, so that the phases of the releases can be shown in Grafana.

\paragraph{Counting every request once.} Two settings keep the proxies from being scraped twice. By default, Istio adds scrape annotations to every pod with a sidecar, so that the generic pod job would scrape it a second time; this merging is switched off in the mesh configuration (\texttt{enablePrometheusMerge: false}), since the \texttt{envoy-stats} job already covers every proxy. The gateway's Helm chart also annotates its pods for scraping, which, before it was disabled by an annotation in the gateway's values, doubled every plain sum of gateway traffic. Ratios such as the error share were not affected, but the request counts of the experiments would have been.

\paragraph{Dashboards that render the same twice.} Argo CD renders every Helm chart again on every sync and compares the result with the cluster. Two charts generated a random value at render time: Grafana its administrator password and Kiali the key that signs its session tokens. Each new value made the Application permanently out of sync, and for Kiali it restarted the pod on every push to the repository. Both values now come from Secrets created by the bootstrap from SOPS, and the charts only refer to them, which I checked by rendering each chart twice and comparing the outputs. Grafana's data source and dashboards (Istio's standard dashboards and a dashboard for rollbacks) are provisioned from files in the repository, so Grafana needs no persistent volume. I had to raise its memory limit twice: 256\,MiB was too little for its first start, which runs the database migrations, and with a limit of 512\,MiB it was killed for running out of memory once dashboards were opened. It now requests 384\,MiB, its idle use, with a limit of 1\,GiB.

%----------------------------------------------------------------------------
\section{Demo application and release pipeline}
\label{sec:ImplDemo}
%----------------------------------------------------------------------------
The demo application is a small HTTP server written in Python with the standard library only. \texttt{GET /api/work} is the endpoint under test, and the only one affected by the faults; \texttt{/health} serves the readiness probe; \texttt{/stats} shows the configuration and counters of the running version, and \texttt{/metrics} its request counters. The faults of \autoref{tab:EvalFaults} are read from environment variables at start-up. The platform first ran podinfo, a common demo application, whose only fault was a random error. I replaced it with my own small application, which already had a CI pipeline on Gitea and could be extended with exactly the faults the experiments needed: errors, fixed and random latency, a failing readiness probe, a memory leak and, later, a wrong answer with a successful status code.

A load generator in the namespace of the application, a \texttt{curl} loop without a sidecar, calls \texttt{/api/work} through the ingress gateway every 0.2\,s. Its requests take the same path as a user's, so the checks always have traffic to measure, at the same rate in every run.

The CI workflow of the application repository runs on Gitea's runner. It builds the image, pushes it to Gitea's registry, clones the platform repository with an access token, rewrites the line \texttt{newTag} of the application's Kustomize file, and commits and pushes the change. Two details had to be adapted. The rewrite uses the POSIX character class \texttt{[[:space:]]} instead of \texttt{\textbackslash s}, which the \texttt{sed} of macOS does not understand, so that the same command can be tested on the control machine and run on the Linux runner. And since CI commits to \texttt{main} while people and the trial runner commit to it too, the platform repository rebases on pull (\texttt{pull.rebase true}), which keeps the history linear.

Since the move to configuration faults, a release under test is the patch shown in \autoref{lst:ImplRelease}, applied to the Rollout by the base Kustomization. The trial runner rewrites the values of this file, commits it and asks Argo CD to refresh.

\begin{lstlisting}[caption={The release file of the demo application (healthy values).},label=lst:ImplRelease]
- op: add
  path: /spec/template/spec/containers/0/env
  value:
    - { name: APP_VERSION, value: "reset" }
    - { name: ERROR_RATE, value: "0" }
    - { name: LATENCY_MS, value: "0" }
    - { name: LATENCY_JITTER_MS, value: "0" }
    - { name: FAIL_READINESS, value: "false" }
    - { name: LEAK_MB_PER_MIN, value: "0" }
    - { name: WRONG_ANSWER, value: "false" }
\end{lstlisting}

%----------------------------------------------------------------------------
\section{Implementing the detection mechanisms}
\label{sec:ImplMechanisms}
%----------------------------------------------------------------------------
I added the mechanisms one at a time and verified each with a healthy and a faulty release before starting the next one. This section describes the points where the design could not be expressed directly in the resources of Argo Rollouts, or where the first version behaved differently from the design.

%----------------------------------------------------------------------------
\subsection{Canary analysis}
\label{sec:ImplCanary}
%----------------------------------------------------------------------------
\paragraph{Errors instead of results at the start.} The first version of the analysis measured immediately when the canary received traffic. At that moment Prometheus had no samples of the canary yet, a rate needs at least two, and the first measurements returned an empty result, which Argo Rollouts records as an \emph{Error}. Errors do not count towards the failure limit, so they did not abort the release, but they were measurements without information. I rejected counting no data as success, by adding \texttt{or vector(1)} to the query, because a canary that receives no traffic at all would then pass. The fix was the 30\,s initial delay of the design, together with a one-minute instead of a two-minute window, which reacts faster and still holds enough requests.

\paragraph{A template change does not restart an aborted release.} After the memory check had aborted a healthy release (\autoref{sec:DesignCanary}), correcting the AnalysisTemplate did not start the release again: Argo Rollouts starts a revision only when the pod template changes, and the template of the analysis is not part of it. I had to retry the release explicitly with \texttt{kubectl argo rollouts retry}. In the same way, a commit that changes only the steps of the strategy starts no release; the new steps apply from the next release on. The experiments rely on this when they switch between strategies (\autoref{sec:EvalStrategies}).

\paragraph{The production overlay replaces the step list.} The production overlay replaces the whole list of steps with its own, longer pauses. When steps were added to the base, they never reached production, and worse, the step index at which the background analysis starts, which is inherited from the base, pointed past the end of production's shorter list, so the analysis would never have started there and every production release would have been promoted unchecked, without any error. Each change of the base steps therefore has a matching change in the production overlay, and I check both renders with \texttt{kubectl kustomize} before a commit.

%----------------------------------------------------------------------------
\subsection{Mirroring}
\label{sec:ImplMirror}
%----------------------------------------------------------------------------
I built mirroring in stages and checked each stage on a healthy release before a faulty one. Three problems appeared that the design had not anticipated.

\paragraph{No pod to mirror to.} In the first attempt the mirror route kept disappearing. The canary's ReplicaSet had zero pods: with traffic routing, Argo Rollouts sizes the canary by its weight, which is 0 during mirroring, and it removes a mirror route that has no destination. The explicit \texttt{setCanaryScale} steps of the design were added for this reason~\cite{ArgoRolloutsCanary}.

\paragraph{The managed route shifts the indexes.} The mirror route is inserted at the top of the VirtualService's route list, so the weighted route moves from index 0 to index 1. With the original exception for \texttt{/spec/http/0} only, Argo CD would have treated the weights at index 1 as drift and reset them in the middle of a release. I widened the exception to the whole route list (\autoref{sec:ImplGitOps}).

\paragraph{The gateway does not see the copies.} The canary analysis selects the canary by the Service that the gateway routed to. During mirroring, however, the gateway reported no requests to \texttt{demo-canary} at all, although the canary was serving copies: Envoy sends the copies without recording them as requests of the caller. I therefore wrote the mirror analysis against the canary pod's own proxy, selected by its pod-template hash. The first faulty release confirmed the design. With 30\,\% of the requests failing, the canary pod's proxy recorded about 2.2 successful and 0.85 failed requests per second, the two measurements of \texttt{mirror-health} read success rates of 0.76 and 0.72, and the release was aborted before any user was routed to the new version. During the same interval, the gateway recorded no 5xx response to users, and after the abort the VirtualService contained neither a mirror route nor a non-zero canary weight.

%----------------------------------------------------------------------------
\subsection{Experiment}
\label{sec:ImplExperiment}
%----------------------------------------------------------------------------
The experiment step copies the pod templates of the stable and the new ReplicaSet (\texttt{specRef: stable} and \texttt{specRef: canary}), so the release under test, whether an image or a configuration, drives it like every other step. Given a weight and a Service name for each template, the controller creates the two Services and adds them to the weighted route of the VirtualService, which then splits the users between stable, baseline and canary. Two details of this setup matter for the checks. The baseline pod has the same pod template as the stable version, so it must not also be selected by the stable Service, or it would receive the stable version's share of the users as well. And the memory check selects the experiment's pods by name, so its patterns must match the names the experiment gives them, but not those of the normal canary pods. The first release ran in dry-run mode and showed values close to zero for a healthy version: an error difference of 0, a latency difference between $-1$ and $+6$\,ms, and no difference in memory growth. Only then were the checks made enforcing.

In the first experiment matrix, the default Rollout ran the experiment at 20\,\% per pod between mirroring and the canary steps. Since the second matrix, the default Rollout has no experiment step (\autoref{sec:DesignExperiment}); the experiment runs at 10\,\% per pod only in the strategy component \texttt{s3-experiment}. Removing the step shifted the index of every later step, so the step at which the background analysis starts had to be changed with it, in the base and in the production overlay (\autoref{sec:ImplCanary}).

%----------------------------------------------------------------------------
\subsection{Blue-green, the smoke test and the load Job}
\label{sec:ImplBlueGreen}
%----------------------------------------------------------------------------
Before relying on the smoke-test Job, I checked the sidecar of a Job pod: in this cluster, Istio injects \texttt{istio-proxy} as an init container with \texttt{restartPolicy: Always}, the native sidecar of Kubernetes~\cite{K8sSidecar}, and the Job completed normally after its last request. The RBAC rules of the Argo Rollouts controller already allowed it to create Jobs.

The smoke test did not appear in Kiali's graph, although its 50 requests had reached the preview pods. Kiali draws its graph from the same Prometheus metrics, and the Job lived for only about 15\,s, shorter than two scrape intervals: its own proxy was never scraped, and the edge recorded by the receiving proxy dropped out of Kiali's default one-minute window soon after the Job ended. Prometheus showed the requests at the preview's pods, all of them at the new version and none at the old one, as the blue-green design intends. Such short-lived pods are one reason why the trial runner counts requests at the receiving proxies (\autoref{sec:EvalMeasurements}).

For the second experiment matrix, the smoke test was replaced as the pre-promotion analysis by \texttt{preview-load}. Its Job runs a shell loop in a \texttt{curl} image that calls the preview Service, waits 0.2\,s and repeats for 60\,s, and it makes a single attempt (\texttt{backoffLimit: 0}) with a hard limit of 90\,s (\texttt{activeDeadlineSeconds}). It ignores the answers and always succeeds; the two Prometheus metrics of the same AnalysisTemplate decide. Argo Rollouts starts all metrics of an AnalysisRun at once, so the Job and the metrics run in parallel, and the 30\,s initial delay of the metrics gives the Job time to start and Prometheus time to scrape its first requests. The first version sent load for 90\,s and made a healthy release exceed the progress deadline, which in blue-green includes the pre-promotion analysis (\autoref{sec:EvalV2Validation}). The strategies of the blue-green application are Kustomize components like those of the canary: \texttt{b0-none} removes both analyses, \texttt{b1-post} the pre-promotion one, \texttt{b2-pre} the post-promotion one, and \texttt{b3-pre-post}, the default, is empty. The smoke test's AnalysisTemplate stays in the base, unused.

%----------------------------------------------------------------------------
\section{Production environment}
\label{sec:ImplProd}
%----------------------------------------------------------------------------
The production environment exists as code but has not been built. Its inventory, its encrypted secrets with its own root CA, and its Argo CD Applications, which track the branch \texttt{release}, are in the repository, and the overlays differ from development as designed: the canary pauses last five and ten minutes; the mirror phase lasts two minutes longer before its check; and Prometheus keeps seven days of data on an 8\,GiB volume instead of 24 hours on 2\,GiB. Building it requires running the same playbook against prgdlv-auto and adding a NetBird route for its pod network (\autoref{sec:ImplNetworking}).

Gitea, which already runs on that VM, could then be served at \texttt{https://gitea.prod.internal} through the production gateway, with a Service without selector and an EndpointSlice pointing at port 3000 of the node. Argo CD would keep addressing Gitea by its IP address, so that development does not come to depend on the production gateway and on name resolution inside pods. Whether production is built is decided at the end of the work (\autoref{sec:DesignEnvironments}); until then, requirement R8 is met by the code but not demonstrated by a running environment.

% Suggested screenshots (guideline 18): the Argo Rollouts dashboard during an
% aborted release, the Grafana rollback dashboard around an abort, and the
% Kiali graph during a canary step. Add them under figures/ when available.

